\documentclass[11pt,a4paper]{article}
\usepackage[T1]{fontenc}
\usepackage[utf8]{inputenc}
\usepackage{lmodern,microtype}
\usepackage[margin=27mm,headheight=15pt]{geometry}
\usepackage{amsmath,amssymb,amsthm,mathtools}
\usepackage{booktabs,tabularx,array,longtable}
\usepackage[table]{xcolor}
\usepackage{enumitem}
\usepackage{fancyhdr,titlesec}
\usepackage{listings}
\usepackage{xurl}
\usepackage[colorlinks=true,linkcolor=navy,citecolor=navy,urlcolor=teal]{hyperref}
\definecolor{navy}{HTML}{183A56}
\definecolor{teal}{HTML}{167D8D}
\definecolor{pale}{HTML}{F1F5F8}
\definecolor{graytext}{HTML}{53606A}
\titleformat{\section}{\Large\bfseries\color{navy}}{\thesection}{.7em}{}
\titleformat{\subsection}{\large\bfseries\color{navy}}{\thesubsection}{.7em}{}
\setlist{topsep=4pt,itemsep=3pt,parsep=1pt}
\pagestyle{fancy}
\fancyhf{}
\fancyhead[L]{\small\color{graytext}Constructive reflection beyond Contender 5}
\fancyhead[R]{\small\color{graytext}Lean 4.19.0}
\fancyfoot[C]{\thepage}
\renewcommand{\headrulewidth}{.3pt}
\newtheorem{theorem}{Theorem}[section]
\newtheorem{lemma}[theorem]{Lemma}
\newtheorem{proposition}[theorem]{Proposition}
\newtheorem{corollary}[theorem]{Corollary}
\theoremstyle{definition}
\newtheorem{definition}[theorem]{Definition}
\newtheorem{remark}[theorem]{Remark}
\newcommand{\N}{\mathbb N}
\newcommand{\llbracket}{\mathopen{[\![}}
\newcommand{\rrbracket}{\mathclose{]\!]}}
\newcommand{\B}{\mathsf B}
\newcommand{\Ty}{\mathsf{Ty}}
\newcommand{\ev}{\operatorname{eval}}
\newcommand{\dep}{\operatorname{depth}}
\newcommand{\pad}{\operatorname{pad}}
\newcommand{\R}{\mathcal R}
\newcommand{\F}{\mathcal F}
\newcommand{\iter}{\operatorname{it}}
\newcommand{\code}[1]{\texttt{\detokenize{#1}}}
\DeclareMathOperator{\cl}{cl}
\DeclareMathOperator{\der}{der}
\lstset{basicstyle=\ttfamily\small,columns=fullflexible,keepspaces=true,
  breaklines=true,frame=single,rulecolor=\color{navy!20},backgroundcolor=\color{pale},
  showstringspaces=false,aboveskip=10pt,belowskip=10pt}
\hypersetup{pdftitle={Constructive Reflection Beyond Contender 5},
  pdfauthor={Research prepared for Vladimir Reshetnikov},
  pdfsubject={An axiom-free Lean large-number entry, exact comparison, and tree-indexed reflection}}
\title{\vspace{-12mm}\textbf{\color{navy}Constructive Reflection\\Beyond Contender 5}\\[6pt]
  \large An axiom-free Lean entry with accessible forests,\\
  structural tree recursion, and unbounded finite-type uniformization}
\author{Research prepared for Vladimir Reshetnikov}
\date{12 October 2026 \quad\textbar\quad Lean 4.19.0}

\begin{document}
\maketitle
\begin{abstract}
We construct a closed, executable Lean natural number and prove that it is
strictly larger than Contender 5 in Cody Roux's \emph{name-the-biggest-number}
repository. The comparison uses a source-faithful Lean port of the original
total raw-term interpreter, including its treatment of ill-typed terms. The
new construction has three substantive ingredients: bounded evaluation in a
language with the full recursor for countably branching trees; a reflection
operator with an accessible countable family of function oracles; and a
structurally generated tree schedule that diagonalizes over an unbounded
dependent family of finite functional types. Every numerical definition is
constructive. The delivered comparison, enumeration, and diagonalization
certificates have empty Lean axiom dependency lists.

The strength statements are explicit. Reflection eventually exceeds every
fixed program in its specified oracle language; successor stages escape the
preceding stage's language; accessible forest stages preserve the exact
values of all their child oracles. We also give finite-path growth theorems,
an independently checked earlier contender that the final entry exceeds,
and a precise explanation of why the construction neither invokes a Busy
Beaver oracle nor defines a universal evaluator for all Lean programs. We
make no claim of an optimal entry under an unspecified machine budget, and
we do not assign an unproved named ordinal to the raw tree schedule.
\end{abstract}

\vspace{2mm}
\noindent\colorbox{pale}{\parbox{\dimexpr\linewidth-2\fboxsep\relax}{
\textbf{Deliverable.} A complete Lean source package, a compiled article,
reproducible checks, original-source provenance, and the theorem
\(\code{Original.contender5}<\code{contender_6}\).
The final numeral is specified and compared, not expanded in decimal.}}

\newpage
\tableofcontents
\newpage

\section{The result and the adapted contest}
\label{sec:rules}

The entry is a natural number \(N\), given by a closed Lean definition named
\code{contender_6}. Its principal certificate is
\[
                 \boxed{C_5<N.}
\]
Here \(C_5\) is the exact numerical interpretation of the repository's fifth
contender, ported to Lean as \code{Original.contender5}. The word ``exact''
refers to the computational equations: the same raw syntax, enumeration,
reverse-index lookup, paired casts, default values, and maximum-selection
algorithm are retained. Section~\ref{sec:baseline} states the cross-language
verification boundary precisely.

The repository asks for constructive definitions, no added axioms, and a
formal strict comparison with the previous contender. It specifies
15 seconds for checking a definition and one minute for its comparison
proof, on a reasonable machine. It discourages defining an entry by applying
a trivial operation to the previous contender's named constant
\cite{contest}. We retain those intentions and substitute Lean for Coq.

\begin{center}
\begin{tabularx}{\linewidth}{@{}p{.24\linewidth}>{\raggedright\arraybackslash}X@{}}
\toprule
Requirement & Lean version used here\\
\midrule
Underlying object & A closed term of type \code{Nat}.\\
Constructivity & Ordinary total definitions and structural recursion; no
\code{noncomputable}, unbounded halting oracle, or partial numerical definition.\\
Axiom policy & The entry and the comparison theorem must each print an empty
axiom dependency list. The supplied audit also checks the supporting theorems.\\
Checking budget & Retain the original 15-second and 60-second targets;
report observed times and the pinned toolchain.\\
Previous entry & Compare with a faithful port of the repository's value,
including malformed raw programs.\\
Use of prior work & Reuse the general interpreter and a semantic embedding,
not the named \code{contender5} constant in the new numerical definitions.\\
Submission & A reviewable archive, without publishing a pull request on the
user's behalf.\\
\bottomrule
\end{tabularx}
\end{center}

The README gives checking deadlines, not a separately specified research or
typing deadline. Thus ``the largest number in the allotted time'' is treated
as a request to construct a strong, fully checked entry during this work.
It does not define a mathematical optimization problem with a fixed machine,
source encoding, and admissible set of programs. With all those parameters
fixed and a hard total checking-time cap, only finitely many source strings
may be readable and an admissible maximum could exist. We make no claim to
have found it. The familiar observation that \(m+1>m\) establishes the absence
of a largest natural number; it does not by itself establish the absence of
a maximum under a fixed finite resource cap.

The useful research question is consequently structural: can a short,
transparent, total definition make an expressive-strength advance over the
existing entry, with a comparison certificate that stays small? The answer
here is yes. The final construction improves a fully verified intermediate
entry and exposes an entire forest of previous evaluators to the next
object language. This latter step is more informative than merely making a
larger finite maximum of their numerical outputs.

\section{Auditing the exact baseline}
\label{sec:baseline}

The inspected repository revision is
\begin{center}
\code{d56dce013e6264292533efdc5e1de71863c6fb86}.
\end{center}
The archive retains the upstream files and their license. The computational
part of \code{Contender.v} is translated in \code{Original.lean}; the original
comparison with Contender 4 is not needed to prove the new result.

\subsection{Simple types and the total raw-term interpreter}

The original type language is
\[
 \sigma,\tau ::= \N\mid \sigma\to\tau,
 \qquad
 \llbracket\N\rrbracket=\N,
 \quad
 \llbracket\sigma\to\tau\rrbracket
    =\llbracket\sigma\rrbracket\to\llbracket\tau\rrbracket.
\]
Terms have variables, annotated lambdas, application, zero, successor, and a
natural-number recursor at each simple result type. Environments contain
dependent pairs \((\sigma,v)\) with
\(v:\llbracket\sigma\rrbracket\). A lambda returns a native total Lean
function. The interpreter is structural recursion on syntax, even though
applying one of the resulting higher-type values can initiate a very long
computation.

Two details matter for an exact comparison. First, the repository's variables
use levels measured from the old end of the environment. If the environment
has length \(\ell\), a variable \(i\) looks up the list position
\(\ell-(i+1)\) when \(i<\ell\), and returns an error value otherwise.
Prepending a newly bound value therefore does not change the meaning of the
old variable levels. Second, casts between unequal types return recursively
defined defaults. At a function type, the default is a constant function
returning the codomain default. The paired arrow casts convert arguments in
one direction and results in the other. These conventions are preserved.

In particular, the enumeration is \emph{not} filtered by an external typing
judgment. All raw terms are interpreted. At top level, a numerical result is
returned as a natural number, while a function result contributes zero. A
proof about a maximum over well-typed closed terms alone would not identify
the target in this repository. The Lean port avoids that substitution.

\subsection{Depth and maximum selection}

Let \(d_T\) be type depth and \(d_0\) original term depth. They satisfy
\begin{align*}
 d_T(\N)&=1,&
 d_T(\sigma\to\tau)&=1+\max(d_T(\sigma),d_T(\tau)),\\
 d_0(\mathsf{var}\ i)&=i+2,&
 d_0(\mathsf{app}\ u\ v)&=1+\max(d_0(u),d_0(v)),\\
 d_0(\lambda_{\sigma,\tau}u)
   &=1+\max(d_T(\sigma),d_T(\tau),d_0(u)).
\end{align*}
Zero and successor have depth one, and a recursor annotated by \(\rho\)
has depth \(1+d_T(\rho)\). Let \(L_0(d)\) be the exact list generated by the
source at bound \(d\), and let \(e_0(t)\) be its numerical evaluation.

The source selects a maximizing term with an accumulator \code{maxBy}. For
proofs it is convenient to use the value-only accumulator
\[
 \operatorname{foldMax}(g,c,[]) = c,\qquad
 \operatorname{foldMax}(g,c,x::xs)
    =\operatorname{foldMax}(g,\max(c,g(x)),xs).
\]

\begin{lemma}[Maximum-selection invariant]
If \(g(b)=c\), then the source's maximum-selection routine returns a term
whose value is \(\operatorname{foldMax}(g,c,xs)\).
\end{lemma}
\begin{proof}
Induct on the list. The empty case is the accumulator invariant. In the
step case, if \(c<g(x)\), both procedures replace their numerical accumulator
by \(g(x)\), and the term accumulator becomes \(x\). Otherwise both numerical
accumulators stay \(c\), and the current best term is retained. The induction
hypothesis applies in either branch.
\end{proof}

Consequently, with the empty maximum equal to zero,
\[
 M_0(d)=\max\bigl(\{0\}\cup\{e_0(t):t\in L_0(d)\}\bigr),
 \qquad C_5=M_0(42).
\]
This equality is proved in Lean, rather than replacing the source definition
silently. Both directions of
\(t\in L_0(d)\iff d_0(t)\le d\) are also checked.

\subsection{What the port certificate does and does not mean}

Lean checks all subsequent comparisons against the translated definition.
The correspondence with the Coq source is a constructor-by-constructor
source audit, not a theorem produced by a verified Coq-to-Lean compiler.
That boundary is recorded in the archive. It is made reviewable by keeping
the original source, pinning its revision, retaining \code{maxBy}, and testing
the delicate lookup, binding, casting, and malformed-term cases. The small
values \(M_0(d)\) for \(d=0,1,2,3,4\) are respectively
\(0,0,1,2,3\). No attempt is made to compute \(M_0(42)\).

\section{A general bounded reflection principle}
\label{sec:reflection}

The central method does not require an ordinal analysis. It needs a total
interpreter, a finite syntax enumeration, and an explicit running input.
The running input is essential for the strongest conclusion: a maximum over
closed descriptions measures description size, whereas a language with an
input constant supports an eventual comparison between functions.

Fix a family of oracles \(\mathcal O\), a numerical interpreter
\(e_{\mathcal O,n}(t)\), and finite lists \(L(d)\) satisfying completeness:
\[
          \dep(t)\le d\quad\Longrightarrow\quad t\in L(d).
\]
For an explicit positive overhead \(p(n)\), consider
\[
 B_{\mathcal O}(n)=p(n)+
     \max\bigl(A_{\mathcal O}(n),
                \max_{t\in L(n+3)}e_{\mathcal O,n}(t)\bigr),
 \qquad p(n)\ge1.
\]
The auxiliary term \(A\) can provide useful uniform inequalities; it does
not compromise totality.

\begin{theorem}[Bounded and eventual diagonalization]
For each term \(t\),
\[
 \dep(t)\le n+3\quad\Longrightarrow\quad
 e_{\mathcal O,n}(t)<B_{\mathcal O}(n).
\]
In particular, the strict inequality holds for every \(n\ge\dep(t)\).
No fixed term \(t\) represents the whole function \(B_{\mathcal O}\) in this
same oracle language.
\end{theorem}
\begin{proof}
Completeness places \(t\) in the finite maximum. Its value is therefore no
larger than that maximum, and adding the positive overhead makes the result
strictly larger. For nonrepresentability, suppose that
\(e_{\mathcal O,n}(t)=B_{\mathcal O}(n)\) for every \(n\), and use
\(n=\dep(t)\). This yields a number strictly smaller than itself.
\end{proof}

This theorem is stronger than a comparison with one selected previous
integer. It supplies a uniform escape from every fixed program of a specified
earlier language. It is also sharply scoped: ``every fixed program'' means
a single finite expression interpreted in that language, not every Lean
function. The meta-level interpreter uses dependent types to interpret all
simple types uniformly. The object language does not thereby acquire a
total evaluator for its own entire meta-language.

The intermediate \code{Reflection.lean} module implements this principle for
the original natural-number language with a unary oracle and an input
constant. Its purpose is to isolate and check the idea before adding tree
recursion. The final entry uses the richer versions below.

\section{Reflection with full countably branching tree recursion}
\label{sec:arboreal}

\subsection{The inductive tree type}

Introduce the strictly positive inductive type
\[
 \B ::= \mathsf Z\mid\mathsf S(a)\mid
          \mathsf L(f),\qquad a:\B,\quad f:\N\to\B.
\]
These are raw well-founded, countably branching trees. A closed tree is
described by a finite program that constructs its branches on demand. It
need not be a finite combinatorial tree. Lean accepts structural recursion
because every \(f(i)\) is a recursive child of \(\mathsf L(f)\).

The recursor with result type \(X\) has parameters
\[
 z:X,\quad s:\B\to X\to X,\quad
 \ell:(\N\to\B)\to(\N\to X)\to X
\]
and equations
\begin{align*}
 \operatorname{rec}_{\B}(z,s,\ell,\mathsf Z)&=z,\\
 \operatorname{rec}_{\B}(z,s,\ell,\mathsf S(a))
   &=s\bigl(a,\operatorname{rec}_{\B}(z,s,\ell,a)\bigr),\\
 \operatorname{rec}_{\B}(z,s,\ell,\mathsf L(f))
   &=\ell\bigl(f,i\mapsto\operatorname{rec}_{\B}(z,s,\ell,f(i))\bigr).
\end{align*}
Thus both the original subtree and its recursive result are available in a
successor branch, and both whole families are available in a limit branch.
This is the full primitive recursor, not just a scalar traversal.

\subsection{The object language \texorpdfstring{\(\mathcal W_f\)}{W(f)}}

Its types are generated by \(\N\), \(\B\), and arrows. In addition to variables,
annotated lambdas, and application, the language contains the following
constants. The result type \(\rho\) ranges over \emph{all} its finite types.

\begin{center}
\begin{tabularx}{\linewidth}{@{}lX@{}}
\toprule
Constant & Type or meaning\\
\midrule
\(0,\mathsf{succ}\) & \(\N\), \(\N\to\N\).\\
\(\mathsf{nrec}_{\rho}\) &
\(\rho\to(\N\to\rho\to\rho)\to\N\to\rho\).\\
\(\mathsf Z,\mathsf S,\mathsf L\) &
\(\B\), \(\B\to\B\), \((\N\to\B)\to\B\).\\
\(\mathsf{brec}_{\rho}\) &
\(\rho\to(\B\to\rho\to\rho)\to
((\N\to\B)\to(\N\to\rho)\to\rho)\to\B\to\rho\).\\
\(\mathsf O\), \(\mathsf X\) &
\(\N\to\N\), \(\N\), interpreted by \(f\) and the current input \(n\).\\
\(\mathsf{old}(t)\) & A natural-number constant with value \(e_0(t)\), for
any raw original term \(t\).\\
\bottomrule
\end{tabularx}
\end{center}

The \(\mathsf{old}\) constructor is an explicit closed numerical embedding.
It does not reinterpret the old term under the new environment. It therefore
preserves exactly the observable value relevant to the contest, including
all prescribed malformed-term results; nonnumeric top-level results
contribute zero. A malformed term need not contribute zero: applying
successor to successor contributes one under the source's error-cast rules.
This avoids an unnecessary identification between the two different
semantic type families.
The new language independently has all the usual simple-type constructors
and its own variables and recursors.

As in the baseline, every raw term has a total interpretation. The default
at \(\B\) is \(\mathsf Z\); other defaults and paired arrow casts are defined
structurally. At top level, a tree or function result contributes zero. This
is a semantics of all raw expressions, not an assertion that all expressions
are conventionally well typed.

Let \(L_W(d)\) be the finite enumeration and let \(e^W_{f,n}\) be its
numerical interpreter. The tree primitives, oracle, and input have depth
one. A recursor annotated by \(\rho\) has depth \(1+d_T(\rho)\), and
\(\dep(\mathsf{old}(t))=d_0(t)\). The remaining depth rules are the familiar
maximum rules for applications and annotated lambdas.

\begin{theorem}[Exact enumeration and totality]
\[
             t\in L_W(d)\quad\Longleftrightarrow\quad\dep(t)\le d.
\]
The evaluator, enumeration, and numerical maxima are total constructive
definitions.
\end{theorem}
\begin{proof}
The forward and backward implications are simultaneous elementary inductions
on the depth bound and the appropriate syntax constructors. A lambda uses
two types and a smaller body; an application uses two smaller terms; a
recursor uses a smaller annotation. The old branch uses the previously
proved original completeness and soundness. The finite lists are built by
finite products, maps, and concatenations.

The interpreter recurses structurally on syntax and on type annotations for
casts. Natural and tree recursors are already total recursors of inductive
types. In particular, evaluation at a function type produces an actual total
function, and a tree-recursion branch receives a family of already defined
recursive values. No search for a normalization certificate is performed.
\end{proof}

\subsection{The unary tree-language reflector}

Define
\begin{equation}
 \R(f)(n)=n+1+\max\left(f(n),\max_{t\in L_W(n+3)}e^W_{f,n}(t)\right).
 \label{eq:R}
\end{equation}
It is \code{Arboreal.Rw} in the source.

\begin{theorem}[Uniform lower bounds and language escape]
For every \(f,n\),
\[
 f(n)<\R(f)(n),\qquad n<\R(f)(n).
\]
Whenever \(d\le n+3\),
\[
 M_0(d)<\R(f)(n).
\]
For every fixed \(t\) in \(\mathcal W_f\),
\[
 n\ge\dep(t)\quad\Longrightarrow\quad
 e^W_{f,n}(t)<\R(f)(n).
\]
In particular \(C_5<\R(f)(42)\), and no fixed \(\mathcal W_f\) expression
represents \(\R(f)\).
\end{theorem}
\begin{proof}
The first two inequalities follow directly from~\eqref{eq:R}. Each old term
of depth at most \(d\) occurs as \(\mathsf{old}(t)\) in the larger enumeration,
with exactly the same numerical value. Thus the whole old maximum is bounded
by the new maximum. The positive overhead makes the inequality strict. The
last statements follow from Section~\ref{sec:reflection}.
\end{proof}

\section{Why a maximum of child values is not an accessible forest}
\label{sec:forest}

A pointwise maximum can majorize a collection of functions while concealing
the individual values needed to execute their programs. For example, knowing
\(\max(f(n),g(n))\) need not reveal whether \(f(n)=0\). An object program can
branch on such a test. We therefore distinguish two operations carefully.

The first hierarchy below uses the finite join
\(\max_{i\le n}f_i(n)\). It certifies eventual numeric domination of each
fixed child, but that alone does not establish that the next stage can
simulate every child-oracle program. The final construction additionally
uses \emph{accessible forests}: a new language receives the full function
\(\Phi:\N\to\N\to\N\), so it can ask for the exact value \(\Phi(i,m)\).

\subsection{Forest syntax and semantics}

The forest language has a base constructor containing any term of
\(\mathcal W\), interpreted with unary oracle \(\Phi(0,\cdot)\) and with the
current input. It also has indexed oracle constants
\(\mathsf O_i:\N\to\N\), a family constant
\(\mathsf{Family}:\N\to\N\to\N\), applications, and annotated lambdas. The
indexed oracle denotes \(\Phi(i,\cdot)\); the family constant denotes \(\Phi\)
itself. This last primitive permits the selected child index to depend on
the running input, rather than requiring a fixed index written into syntax.

The base constructor preserves the full environment. Consequently the
forest language retains the tree and natural-number recursors at every
simple result type, and can combine them with calls to any of its child
oracles. An indexed oracle has depth \(i+1\), the family constant has depth
one, and base terms retain their \(\mathcal W\)-depth. The enumerator is finite
at every bound: its indexed oracle branch ranges over only finitely many
indices. Write its list as \(L_F(d)\) and its numerical semantics as
\(e^F_{\Phi,n}\).

Define the prefix maximum
\[
 P_\Phi(n)=\max_{0\le i\le n}\Phi(i,n)
\]
and the forest reflector
\begin{equation}
 \F(\Phi)(n)=n+2+
      \max\left(P_\Phi(n),
                \max_{t\in L_F(n+3)}e^F_{\Phi,n}(t)\right).
 \label{eq:F}
\end{equation}
The extra unit relative to~\eqref{eq:R} makes a useful comparison strict.

\begin{theorem}[Accessible forest reflection]
For \(i\le n\),
\[
             \Phi(i,n)<\F(\Phi)(n).
\]
For each fixed forest expression \(t\),
\[
 n\ge\dep(t)\quad\Longrightarrow\quad
 e^F_{\Phi,n}(t)<\F(\Phi)(n).
\]
Also
\[
            \R(\Phi(0,\cdot))(n)<\F(\Phi)(n).
\]
No fixed expression in the language with this forest oracle represents the
whole function \(\F(\Phi)\).
\end{theorem}
\begin{proof}
The prefix maximum contains \(\Phi(i,n)\). Completeness contains every
expression of sufficiently small depth. The respective positive overheads
give the first two claims. For the comparison with \(\R\), every
\(\mathcal W\)-term occurs as a forest base term at the same depth, and
\(P_\Phi(n)\ge\Phi(0,n)\). Thus the maximum in~\eqref{eq:F} is at least the
maximum in~\eqref{eq:R}; the overhead is larger by one. Nonrepresentability
again uses the depth of a proposed representing expression as the diagonal
input.
\end{proof}

The theorem quantifies over the actual forest-language semantics. A program
that combines finitely many child functions, or selects a child uniformly
through \(\mathsf{Family}\), is beaten once its finite description fits the
depth budget. This is the precise sense in which the forest construction
repairs the information loss of a numerical maximum. It does not assert a
reflection principle for arbitrary formal theories or for all meta-level
Lean definitions.

The checked escape theorem concerns the forest syntax just specified.
The observation that the primitives permit a syntax-directed translation
of individual child-oracle programs is not presented as a separately
machine-checked generic translation theorem; quantitative verified
translations remain a proposed next step.

\section{A large, explicit structural schedule}
\label{sec:schedule}

Both reflection operators can already produce large integers. We now give
an explicit schedule for iterating them that uses more than a fixed finite
height. All definitions in this section construct raw \(\B\)-trees. Their
correctness as total programs is independent of any assignment of ordinal
names.

\subsection{Tree arithmetic and Veblen-style operations}

For any type \(X\), write
\[
 \operatorname{fold}(z,s,\ell,\mathsf Z)=z,\quad
 \operatorname{fold}(z,s,\ell,\mathsf S(a))
    =s(\operatorname{fold}(z,s,\ell,a)),
\]
with the analogous limit equation
\(\ell(i\mapsto\operatorname{fold}(z,s,\ell,f(i)))\).
Define iteration by
\(\iter(u,0,x)=x\) and
\(\iter(u,k+1,x)=\iter(u,k,u(x))\). Let
\[
 \underline k=\iter(\mathsf S,k,\mathsf Z),\qquad
 \omega_{\B}=\mathsf L(k\mapsto\underline k).
\]
The structural arithmetic operations are
\begin{align*}
 a+b&=\operatorname{fold}(a,\mathsf S,\mathsf L,b),\\
 a\cdot b&=\operatorname{fold}
       (\mathsf Z,r\mapsto r+a,\mathsf L,b),\\
 a^b&=\operatorname{fold}
       (\mathsf S(\mathsf Z),r\mapsto r\cdot a,\mathsf L,b).
\end{align*}
Put \(p(a)=\omega_{\B}^{a}\). For a tree transformation \(u\), define
\begin{align*}
 \cl(u,a)&=\mathsf L(k\mapsto\iter(u,k,\mathsf S(a))),\\
 \der(u)&=\operatorname{fold}
       (\cl(u,\mathsf Z),a\mapsto\cl(u,a),\mathsf L),\\
 \operatorname{common}(u_0,u_1,\ldots)
   &=\der\bigl(a\mapsto\mathsf L(k\mapsto u_k(a))\bigr),\\
 V_u&=\operatorname{fold}(u,\der,\operatorname{common}).
\end{align*}
Thus \(V_u(a)\) is a tree transformation and \(V_p\) is a Veblen-style
schedule. These equations are the specification, not an appeal to an
informally understood large ordinal. Their names explain the construction's
inspiration. Establishing equality with a particular canonical Veblen
notation system would require additional invariants and a separate proof.

\subsection{Uniformization over unbounded finite functional levels}

A more consequential step is to organize all finite levels of higher-order
transformations in a dependent family:
\[
 E_0=\B\to\B,\qquad E_{k+1}=E_k\to E_k.
\]
This is an actual type-valued recursive definition in Lean. At level zero,
use \(\operatorname{common}\) for countable limits and the derivative of a
diagonal evaluation for tree-indexed diagonals. Lift these operations
pointwise:
\begin{align*}
 \Lambda_0(u_i)&=\operatorname{common}(u_i),&
 \Delta_0(v_a)&=\der(a\mapsto v_a(\mathsf Z)),\\
 \Lambda_{k+1}(u_i)(x)&=\Lambda_k(i\mapsto u_i(x)),&
 \Delta_{k+1}(v_a)(x)&=\Delta_k(a\mapsto v_a(x)).
\end{align*}
For \(s:E_{k+1}\) and \(z:E_k\), define
\[
 Q_k(s)(z)=
       \Delta_k\bigl(a\mapsto\operatorname{fold}(z,s,\Lambda_k,a)\bigr).
\]
Then \(Q_k:E_{k+2}\). Choose canonical seeds
\[
 q_0=p,\qquad q_1=\der,\qquad q_{k+2}=Q_k.
\]
The operation \(\operatorname{collapse}_k:E_k\to E_0\) applies a functional
to the canonical seed one level below until level zero is reached:
\[
 \operatorname{collapse}_0(u)=u,\qquad
 \operatorname{collapse}_{k+1}(u)
      =\operatorname{collapse}_k(u(q_k)).
\]
Two useful checks reduce by definition:
\begin{align*}
 \operatorname{collapse}_1(q_1)&=\der(p),\\
 \operatorname{collapse}_2(q_2)
   &=\der\bigl(a\mapsto V_p(a)(\mathsf Z)\bigr).
\end{align*}
These identities are checked by \code{rfl} in Lean.

Finally set
\begin{align}
 U(a)&=\mathsf L\bigl(k\mapsto
   \operatorname{collapse}_{k+2}(q_{k+2})(\mathsf S(a))\bigr),\label{eq:U}\\
 \Theta&=\mathsf L\bigl(m\mapsto\iter(U,m+1,\mathsf Z)\bigr).
 \label{eq:theta}
\end{align}
This \(\Theta\) is \code{treeIndex}. Equation~\eqref{eq:U} ranges over
unbounded finite functional levels in one uniformly typed definition;
equation~\eqref{eq:theta} then iterates that operation. The integer 42 does
not appear in this construction.

Dependent uniformization of this kind is motivated by the distinction
between fixed simple types and a universe containing a type family indexed
by natural numbers. Escard\'o gives an explicit development of related
Church encodings and Brouwer trees \cite{escardo}. We use that perspective
to design a schedule; the delivered definitions and proofs stand on their
own. No exercise postulate from that development is imported.

\section{Two hierarchies and the final entry}
\label{sec:entry}

\subsection{The unary reflection hierarchy}

Define functions \(H_a:\N\to\N\) by structural recursion on \(a:\B\):
\begin{align}
 H_{\mathsf Z}&=\R(\mathsf{succ}),\nonumber\\
 H_{\mathsf S(a)}&=\R(H_a),\label{eq:H}\\
 H_{\mathsf L(c)}&=
       \R\bigl(n\mapsto\max_{i\le n}H_{c(i)}(n)\bigr).\nonumber
\end{align}
Every call in the limit branch is on an immediate recursive child. The
argument \(n\) may vary, but the structurally decreasing argument is the
tree. This is a standard higher-type use of the recursor.

\begin{theorem}[Successors, limits, and bounded paths]
For every \(a,n\), \(H_a(n)<H_{\mathsf S(a)}(n)\). For \(i\le n\),
\[
                    H_{c(i)}(n)<H_{\mathsf L(c)}(n).
\]
If a path of \(r\) upward steps uses successor edges or limit edges whose
selected indices are at most \(n\), then
\[
                    H_a(n)+r\le H_b(n)
\]
for its starting and ending trees \(a,b\).
\end{theorem}
\begin{proof}
Successor strictness is \(f(n)<\R(f)(n)\). At a limit, the selected child is
bounded by the prefix maximum, which is strictly below its reflection.
Induction on the finite path adds one to the certified gap at every step.
No monotonicity assumption on the reflector's oracle argument is used.
\end{proof}

The same outer-constructor argument proves
\[
                         C_5<H_a(42)
                         \quad\hbox{for every }a:\B.
\]
For this statement it is not necessary to traverse the potentially enormous
tree: each of its three possible outer constructors already applies \(\R\).

Define \(\pad(a,0)=a\) and
\(\pad(a,k+1)=\mathsf S(\pad(a,k))\), and put
\begin{equation}
            C=\mathsf L(k\mapsto\pad(\Theta,k+1)),
            \qquad B(n)=H_C(n).
\label{eq:C}
\end{equation}
Padding is literal iteration of \(\R\) at the function level:
\[
             H_{\pad(a,k)}=\R^{[k]}(H_a).
\]
This function equality is established by structural induction and
congruence; function extensionality is unnecessary.

\begin{corollary}[Diagonalization in reflection height]
\[
            \R^{[n+1]}(H_\Theta)(n)<B(n).
\]
In particular \(\R^{[43]}(H_\Theta)(42)<B(42)\), and \(C_5<B(42)\).
\end{corollary}
\begin{proof}
Select child \(i=n\) in~\eqref{eq:C} and use the padding identity.
\end{proof}

\subsection{The accessible forest hierarchy}

Write \(\overline f\) for the constant family \(i\mapsto f\). Define
\begin{align}
 K_{\mathsf Z}&=\F(\overline{\mathsf{succ}}),\nonumber\\
 K_{\mathsf S(a)}&=\F(\overline{K_a}),\label{eq:K}\\
 K_{\mathsf L(c)}&=\F(i\mapsto K_{c(i)}).\nonumber
\end{align}
The crucial distinction from~\eqref{eq:H} is that the limit reflector
receives the exact child family, not just its maximum at the current input.
The family primitive can call any \(K_{c(i)}(m)\) requested by a finite program.

For every \(i\le n\), the strict child inequality
\(K_{c(i)}(n)<K_{\mathsf L(c)}(n)\) follows from forest reflection. At a
successor, the corresponding forest program can use the preceding function
uniformly, and the next stage eventually exceeds every fixed expression in
that language. At a limit, it can also use uniformly selected children.
These claims follow from the same exact finite-depth theorem rather than
from an assumed monotonicity principle.

\subsection{The final oracle forest and natural number}

The final entry strictly improves the completed unary candidate as well as
using the accessible-forest hierarchy. Define the family \(\Psi\) by
\begin{equation}
 \Psi(0,n)=B(n),\qquad
 \Psi(k+1,n)=K_{\pad(C,k)}(n).
 \label{eq:psi}
\end{equation}
Our proposed contest number is
\begin{equation}
                  \boxed{N=\F(\Psi)(42).}
 \label{eq:N}
\end{equation}
Every component in~\eqref{eq:N} has been specified by finite syntax and total
structural recursion. In particular, \(\Psi\) is not a list of arbitrary
unprovided functions. It is one closed program producing the appropriate
function for any requested index.

\begin{theorem}[Final comparison and explicit structural margins]
The final entry satisfies
\begin{align*}
 C_5&<B(42)<N,\\
 \R^{[43]}(H_\Theta)(42)&<B(42)<N,\\
 K_{\pad(C,41)}(42)&<N.
\end{align*}
The last expression is 41 additional successor forest-reflections above
\(K_C\). Moreover \(K_C(42)\) itself strictly exceeds the child with
43 successor forest-reflections above \(K_\Theta\).
\end{theorem}
\begin{proof}
The first chain uses the proved unary candidate inequality and oracle zero
in~\eqref{eq:psi}. The second is its reflection-height lower bound. For the
third, select oracle index 42, for which the parameter \(k\) is 41. It is
included in the prefix maximum because \(42\le42\). The final assertion
selects child 42 in the definition of \(C\) and applies the forest child
inequality.
\end{proof}

The join in~\eqref{eq:psi} is intentional. Although both hierarchies are
strong, a comparison between them cannot be obtained just by substituting
a larger oracle into every raw program. Including the earlier candidate as
an exact accessible oracle gives a short certified improvement without that
invalid assumption. The final entry also reflects over the whole joined
language; it is not merely the maximum of two named integers.

\section{Limits of the strength claims}
\label{sec:limits}

\subsection{Raw trees are not canonical ordinals}

If one interprets \(\mathsf Z\) as zero, \(\mathsf S\) as successor, and
\(\mathsf L\) as a countable supremum, then
\(\mathsf Z\) and \(\mathsf L(i\mapsto\mathsf Z)\) have the same ordinal
value. Nevertheless, the reflected hierarchy satisfies
\[
       H_{\mathsf Z}(n)<H_{\mathsf L(i\mapsto\mathsf Z)}(n).
\]
It therefore does not descend through this ordinal interpretation. This is
an explicit counterexample to rank invariance, not just a missing proof.

Consequently, we do not label \(\Theta\) as \(\Gamma_0\), the small Veblen
ordinal, or the Bachmann--Howard ordinal, and we do not identify \(H\) with a
standard fast-growing hierarchy. The raw construction is fully specified
and its inequalities are proved directly. Type-theoretic ordinal research
distinguishes several tree and order presentations precisely because their
constructive operations and equality principles differ \cite{ordinals}.

\subsection{Oracle order and input order}

Neither
\[
 f\le g\ \Longrightarrow\ \R(f)\le\R(g)
 \qquad\hbox{nor}\qquad
 n\le m\ \Longrightarrow\ \R(f)(n)\le\R(f)(m)
\]
is used. An individual interpreted program can decrease when an oracle
value or input increases. Enlarging the finite set of programs does not
by itself prove either desired inequality when the interpretation of every
program changes at the same time. The delivered proofs use only exact
embeddings, explicit prefix maxima, positivity, and structural recursion.

This distinction also explains why the unary limit in~\eqref{eq:H} has a
weaker language-level interpretation than the forest limit in~\eqref{eq:K}.
The former majorizes each child numerically; the latter provides every child
as a callable function. The paper keeps these statements separate.

\subsection{A function theorem does not rank every fixed integer}

Eventual domination establishes a meaningful separation between functions.
It does not assert that our value at 42 exceeds every value at 42 of every
definable function: there are arbitrarily long finite programs for arbitrarily
large constants. At the selected input, the direct maximum certificate
covers the specified depth bound 45. The inequalities over \(C_5\), the
intermediate candidate, and the displayed reflection iterates are separate,
finite-input theorems.

\subsection{Constructive totality and the absence of Busy Beaver search}

For each finite syntax bound the enumerated list is finite, and every entry
has a total interpreter by construction. We do not enumerate arbitrary
Turing machines and wait to discover which ones halt. No undecidable
predicate selects a subset of programs. No maximum over all Lean definitions
is taken. All higher-type functions used in the closed entry are themselves
given by executable definitions.

The construction can still be computationally infeasible by an immense
margin. Constructive existence does not imply feasible expansion. Even a
very small syntax enumeration grows rapidly. If \(A_d\) counts tree-language
types, \(O_d\) original terms, and \(T_d\) tree-language terms, the exact list
construction gives
\[
 A_0=T_0=0,\qquad A_{d+1}=2+A_d^2,
\]
\[
 T_{d+1}=O_{d+1}+d+A_d^2T_d+T_d^2+7+2A_d.
\]
These counts describe syntax, before any enormous semantic evaluation is
attempted. The comparison proofs use membership and inequalities rather
than materializing those lists at the contest bound.

\section{Formal verification and reproducibility}
\label{sec:verification}

\subsection{Axiom discipline}

Lean's standard ecosystem includes principles such as propositional
extensionality, quotient soundness, and classical choice; it also exposes a
command reporting the axioms on which a declaration depends
\cite{leanaxioms}. The contest's rule is implemented here by requiring an
empty report, including for the main inequality. This is stricter than
merely avoiding a user-written \code{axiom} declaration.

Some convenient library lemmas may carry stronger dependency sets than their
elementary statements suggest. We therefore supply small structural proofs
of the needed maximum and list lemmas. The resulting certificate does not
use function extensionality, classical choice, quotient soundness,
\code{sorryAx}, or compiler-trusting native decision procedures. An empty
axiom list does not prove Lean's kernel soundness inside Lean; the usual
trust in the checker and its foundational rules remains.

The theorem name follows the original contest convention:
\begin{lstlisting}
theorem contender_5_lt_contender_6 :
    Original.contender5 < contender_6 := ...

#print axioms contender_6
#print axioms contender_5_lt_contender_6
\end{lstlisting}
The archive includes the complete proof in place of the schematic ellipsis.
Its audit output is retained separately. The displayed snippet explains
what to inspect; it is not a replacement for the checked file.

\subsection{Why typechecking remains short}

A checker needs to verify that a recursive definition is well formed and
that its recursive calls are justified. It does not normally normalize the
defined natural number. Lean compiles structural recursion into recursors;
the relevant mechanism is documented in its reference manual
\cite{leanrecursion}.

There is, however, an engineering hazard. A proof that rewrites a huge
specialized closed expression can accidentally make the elaborator unfold
the tree schedule or start reducing a semantic computation. The proof
package avoids that hazard by proving the needed result for an abstract tree
and an abstract reflector first, then specializing the theorem. The abstract
proof has no reason to inspect the internal structure of \(\Theta\).

This is ordinary abstraction in a proof, not an opaque axiom. Every helper
definition and proof is included and checked from source. The timed build
regenerates the project's compiled modules in dependency order; only the
standard Lean distribution is reused precompiled.

\subsection{Observed build and regression checks}

The complete source build passed with Lean 4.19.0. Rebuilding all eight
modules in the final entry's dependency graph, including their proofs, took
\textbf{2.962 seconds}. Rebuilding all eleven delivered mathematical modules,
including the companion reflection module, smoke checks, and the full axiom
audit, took \textbf{4.247 seconds} wall-clock time. Both are below the
original 15-second definition target on this host; the entire check is also
below the one-minute proof target. These are observed wall-clock timings,
not hardware-independent guarantees.

\begin{center}
\begin{tabular}{@{}lr@{}}
\toprule
Source module & Seconds\\
\midrule
\code{Original.lean} & 0.706\\
\code{Trees.lean} & 0.262\\
\code{Reflection.lean} & 0.336\\
\code{Hierarchy.lean} & 0.261\\
\code{Arboreal.lean} & 0.884\\
\code{Forest.lean} & 0.381\\
\code{Preliminary.lean} & 0.158\\
\code{ForestHierarchy.lean} & 0.161\\
\code{Contender.lean} & 0.150\\
\code{Smoke.lean} & 0.318\\
\code{Audit.lean} & 0.629\\
\bottomrule
\end{tabular}
\end{center}

All \textbf{288 explicitly declared definitions, inductive types, and theorems}
in the delivered mathematical and smoke-test modules were audited;
\textbf{all 288 axiom dependency lists were empty}. A separate traversal of
the elaborated dependency graph found 502 reachable declarations, including
169 from this package. Neither the original \code{contender5} constant nor
the intermediate named number appears in the final number's definition
closure. There are no axioms, unsafe or partial declarations, or package
implementation overrides in that closure. The ordinary core tags for
\code{Nat.modCore} and its helper, and the core representation optimization
for \code{Nat.repr}, are reported explicitly in the audit output; they are
not package additions or proof-certification assumptions.

The independent Python/Lean regression comparison passed for
\textbf{25,134 original term values} in enumeration order:
\begin{center}
\begin{tabular}{@{}rrrr@{}}
\toprule
Depth & Types & Terms & Maximum\\
\midrule
0 & 0 & 0 & 0\\
1 & 1 & 2 & 0\\
2 & 2 & 10 & 1\\
3 & 5 & 146 & 2\\
4 & 26 & 24,976 & 3\\
\bottomrule
\end{tabular}
\end{center}
The model comparison and the declaration-closure inspection are supplementary
validation. The numerical comparison theorem itself is checked by Lean's
kernel. Neither large contender was evaluated.


The smoke checks exercise the numerical and higher-type semantics rather
than attempting the contest integer. They cover original reverse lookup,
out-of-range variables, binder order, malformed application, mismatching
casts, closed legacy nodes under a new environment, and tree recursion with
natural-number, function, and tree-valued results. The latter includes a
countably branching input. Forest checks also select a child index beyond
the prefix maximum, pass a bound variable to the family primitive, and call
the family from inside successor and limit branches of tree recursion.
These are kernel reduction checks at small
instances, not a substitute for the general theorems.

The complete source-level audit and reproduction instructions are supplied
alongside the article. The recorded runtime uses the Lean 4.19.0 release
kernel and standard libraries. In this execution environment, a small
application-path launcher is needed because the usual executable-path
query is unavailable. Its source is included for transparency; it changes
application-path discovery, not the Lean kernel or elaborator. A normal
Lean 4.19.0 installation uses the provided build commands directly without
this environment-specific launcher.

\section{Further research questions}
\label{sec:research}

The following directions aim at stronger mathematical information or more
economical certified entries. They are proposed research, not claims proved
by this archive.

\begin{enumerate}[label=\textbf{\arabic*.},leftmargin=*]
\item \textbf{Quantitative translations between object languages.}
Construct verified syntax translations from ordinary typed System T,
tree-recursive languages, and their forest extensions, with explicit depth
overhead. A logical relation at every finite type can avoid assuming
function extensionality. Such translations would turn informal comparisons
of language expressivity into exact input thresholds.

\item \textbf{Ordinal calibration on a restricted tree sublanguage.}
Choose canonical trees with increasing cofinal sequences and prove the
arithmetic and fixed-point invariants they need. Determine which restricted
schedules admit a faithful ordinal interpretation. The raw reflection
hierarchy itself cannot simply be quotiented by ordinal rank, as the
counterexample in Section~\ref{sec:limits} shows.

\item \textbf{Lower bounds in established hierarchies.}
Prove explicit comparisons with familiar fast-growing functions or
normalization bounds. The aim should be certified inequalities with stated
fundamental sequences and numerical thresholds, rather than attaching a
large ordinal name to a program.

\item \textbf{A general theorem for accessible oracle forests.}
Abstract the family-language construction over a syntactic category and
prove preservation of exact oracle access under arbitrary finite
composition. The present implementation supplies one concrete total
language and the key diagonal inequality; a general metatheorem could
organize many further contests.

\item \textbf{More expressive, still total object languages.}
Add further strictly positive inductive types or carefully stratified
dependent type formers, each with a verified interpreter. A new language
must remain below its meta-language in the relevant sense; unrestricted
self-interpretation cannot be inserted by assertion.

\item \textbf{Uniform reflection over finite universe fragments.}
Investigate a sequence of syntactic universes with explicit interpretation
maps and verified consistency boundaries. This may give a stronger systematic
advance than adding further successor padding. It is also substantially
harder than the present dependent finite-type uniformization.

\item \textbf{Semantic preservation of normalization optimizations.}
Replace explicit Cartesian-product enumeration by sharing, dynamic
programming, or canonical syntax DAGs, while proving that the numerical
maximum is unchanged. The objective is useful small-instance experimentation
and a cleaner implementation, not feasible evaluation of the final entry.

\item \textbf{An official, reproducible Lean contest protocol.}
Fix the toolchain, machine or normalized benchmark, byte accounting, allowed
imports, axiom policy, and treatment of compiled dependency caches. Separate
definition checking from normalization explicitly. This would make future
claims of admissibility substantially easier to compare.

\item \textbf{A verified source bridge for the old contender.}
Formalize an intermediate syntax/semantics specification that can be checked
from both Coq and Lean, or establish a verified translation for this fragment.
That would remove the manual cross-language source audit from the baseline
identity argument.

\item \textbf{Economy of the final construction.}
Search for a shorter entry with the same or stronger certified lower bounds,
and for stronger entries with comparable check time. Optimizing a precisely
defined language and checker budget is a concrete experimental problem;
there is no evidence here of a globally optimal submission.
\end{enumerate}

\section{Conclusion}

The final definition is a concrete constructive natural number with a checked
strict comparison against the repository's fifth contender. Its main advance
is the combination of total higher-type interpretation, explicit access to
forests of earlier evaluators, and a dependent family of structural schedules.
The exact finite-depth diagonal theorem explains what each reflection step
achieves. The separate fixed-input certificates explain why the chosen entry
wins the required comparison.

The strongest claims in the paper are backed by Lean proofs, while the
remaining boundaries are explicit: source-faithful cross-language porting,
machine-specific timing, raw trees rather than canonical ordinals, and no
universal evaluation of all Lean programs. The result is a reproducible
entry with explicit, checked strength claims.

\clearpage
\appendix
\section{Guide to the source modules}

\begin{longtable}{@{}p{.27\linewidth}p{.67\linewidth}@{}}
\toprule
Module & Role\\
\midrule\endhead
\code{Original.lean} & Exact computational port; original enumeration,
maximum invariant, membership characterizations.\\
\code{Reflection.lean} & Unary reflection for the original natural-number
language; a smaller independently checked diagonalization module.\\
\code{Trees.lean} & Raw tree type, arithmetic, functional type family,
uniformization schedule, and auxiliary fast-tree lemmas.\\
\code{Arboreal.lean} & Full tree-recursive object language and unary
reflector \(\R\).\\
\code{Hierarchy.lean} & Generic tree-indexed unary reflection, prefix
maximum, padding, and path-growth theorems.\\
\code{Forest.lean} & Accessible countable oracle family and forest
reflector \(\F\).\\
\code{Preliminary.lean} & The completed unary candidate family, its
strict source comparison, and its 43-reflection lower bound.\\
\code{ForestHierarchy.lean} & Structural stages with exact access to
their entire child families and corresponding growth certificates.\\
\code{Contender.lean} & The final joined oracle forest, natural number,
and fixed-input comparison certificates.\\
\code{Smoke.lean} & Small exact semantic regression checks.\\
\code{Audit.lean} & Explicit axiom dependency checks of the deliverable.\\
\bottomrule
\end{longtable}

The reproduction scripts and the README state the exact build order. Source
proofs, numerical definitions, and verification logs are all included. The
upstream MIT license is retained with the source-faithful port and the
original files.

\begin{thebibliography}{9}
\bibitem{contest}
Cody Roux and contributors.
\emph{name-the-biggest-number}: README and \code{Contender.v}.
Repository revision \code{d56dce013e6264292533efdc5e1de71863c6fb86}.
\url{https://github.com/codyroux/name-the-biggest-number}.
Accessed 12 October 2026. The original Contender 5 is credited in the source
to Samuel Gr\"utter.

\bibitem{escardo}
Mart\'in Escard\'o.
\emph{Ordinals in Goedel's system T and in Martin-Loef Type Theory}.
2010; Agda development, November 2011.
\url{https://martinescardo.github.io/papers/ordinals/ordinals.html}.

\bibitem{ordinals}
Nicolai Kraus, Fredrik Nordvall Forsberg, and Chuangjie Xu.
\emph{Type-Theoretic Approaches to Ordinals}.
arXiv:2208.03844.
\url{https://arxiv.org/abs/2208.03844}.

\bibitem{leanaxioms}
The Lean community.
\emph{Theorem Proving in Lean 4}, ``Axioms and Computation.''
\url{https://lean-lang.org/theorem_proving_in_lean4/Axioms-and-Computation/}.

\bibitem{leanrecursion}
The Lean community.
\emph{Lean Language Reference}, version 4.19.0,
``Recursive Definitions.''
\url{https://lean-lang.org/doc/reference/4.19.0/Definitions/Recursive-Definitions/}.
\end{thebibliography}

\end{document}
